\documentclass[a4paper]{article}

% Basic packages
% Español (México) con XeLaTeX: fontspec para fuentes Unicode, polyglossia
% para la división silábica y los nombres del documento en la variante mexicana.
\usepackage{fontspec}
\usepackage{polyglossia}
\setdefaultlanguage[variant=mexican]{spanish}
% Los nombres chinos van como «pinyin (original)»: xeCJK da a los caracteres
% Han su propia fuente; sin ella XeLaTeX los omite con un aviso.
% FandolSong no cubre algunos caracteres de nombres (U+73FA); WenQuanYi Zen Hei sí.
\usepackage[AutoFallBack=true]{xeCJK}
\setCJKmainfont{FandolSong-Regular.otf}
\setCJKfallbackfamilyfont{\CJKrmdefault}{WenQuanYi Zen Hei}
% polyglossia no traduce los nombres de listings.
\AtBeginDocument{\providecommand{\lstlistingname}{}\renewcommand{\lstlistingname}{Listado}%
  \providecommand{\lstlistlistingname}{}\renewcommand{\lstlistlistingname}{Índice de listados}}
\usepackage{amsmath, amssymb}
\usepackage{graphicx}
\usepackage[margin=2.5cm]{geometry}
\usepackage[most]{tcolorbox}
\usepackage{etoolbox}
\usepackage{listings}
\usepackage{hyperref}
\usepackage{booktabs}
\usepackage{subcaption}
\usepackage{float}

% Highlight boxes
\newtcolorbox{knowledgebox}[1]{
    enhanced, colback=blue!5!white, colframe=blue!75!black, colbacktitle=blue!75!black,
    coltitle=white, fonttitle=\bfseries, title={#1}, attach boxed title to top left={yshift=-2mm, xshift=2mm},
    boxrule=1pt, sharp corners
}

\newtcolorbox{importantbox}[1]{
    enhanced, colback=yellow!10!white, colframe=yellow!80!black, colbacktitle=yellow!80!black,
    coltitle=black, fonttitle=\bfseries, title={#1}, sharp corners
}

\newtcolorbox{warningbox}[1]{
    enhanced, colback=red!5!white, colframe=red!75!black, colbacktitle=red!75!black,
    coltitle=white, fonttitle=\bfseries, title={#1}, sharp corners
}

% Listings
\lstset{
    language=Python,
    basicstyle=\ttfamily\small,
    keywordstyle=\color{blue},
    stringstyle=\color{red!60!black},
    commentstyle=\color{green!60!black},
    breaklines=true,
    frame=single,
    numbers=left,
    numberstyle=\tiny\color{gray},
    captionpos=b,
    extendedchars=true
}

% Front-page metadata
\newcommand{\notetitle}{MIT 6.S191 Lecture 5: Deep Reinforcement Learning}
\newcommand{\noteauthors}{Elaboradas a partir de la clase de Alexander Amini}
\newcommand{\notedate}{Primavera de 2025}
\newcommand{\videochannel}{Alexander Amini (MIT)}
\newcommand{\videopublishdate}{2025-03-31}
\newcommand{\videoduration}{62 minutos}
\newcommand{\videourl}{https://www.youtube.com/watch?v=to-lHJfK4pw}
\newcommand{\videocoverpath}{cover.jpg}
\newcommand{\repourl}{https://github.com/hqhq1025/ai-course-notes}

% Footer with repo info
\usepackage{fancyhdr}
\pagestyle{fancy}
\fancyhf{}
\fancyfoot[C]{\thepage}
\fancyfoot[R]{\footnotesize\color{gray}\href{\repourl}{AI Course Notes}}
\renewcommand{\headrulewidth}{0pt}
\renewcommand{\footrulewidth}{0pt}

\begin{document}

\begin{titlepage}
\centering
{\Large Notas del curso\par}
\vspace{1.2cm}
{\huge\bfseries \notetitle\par}
\vspace{0.8cm}
{\large \noteauthors\par}
\vspace{0.3cm}
{\large \notedate\par}
\vspace{1.2cm}

\ifdefempty{\videocoverpath}{
{\small Al generar las notas, indique la ruta local de la portada del video.\par}
}{
\includegraphics[width=0.82\textwidth,height=0.45\textheight,keepaspectratio]{\videocoverpath}\par
}

\vfill
\begin{tcolorbox}[width=0.9\textwidth, colback=black!2!white, colframe=black!60, sharp corners]
\textbf{Autor o canal del video}: \videochannel\par
\textbf{Fecha de publicación}: \videopublishdate\par
\textbf{Duración del video}: \videoduration\par
\ifdefempty{\videourl}{}{
\textbf{Enlace del video}: \href{\videourl}{\nolinkurl{\videourl}}\par
}\par
\textbf{Página del proyecto}: \href{\repourl}{\nolinkurl{\repourl}}\par
\end{tcolorbox}
\end{titlepage}

\tableofcontents
\newpage

%% --- Inicio del contenido --- %%

\section{Introducción: aprender en entornos dinámicos}

Esta clase es la quinta de la serie de MIT 6.S191 (Introduction to Deep Learning), con el tema de \textbf{Deep Reinforcement Learning} (aprendizaje por refuerzo profundo). En las cuatro clases anteriores, el curso ya cubrió las arquitecturas centrales del aprendizaje profundo (redes completamente conectadas, RNN, CNN) y dos paradigmas de aprendizaje (Supervised Learning y Unsupervised Learning). Esta clase introduce el tercer paradigma de aprendizaje —\textbf{Reinforcement Learning} (aprendizaje por refuerzo)— y explora cómo combinar las técnicas de aprendizaje profundo con el aprendizaje por refuerzo para que un modelo pueda aprender una política óptima de comportamiento en un \textbf{entorno dinámico} mediante la interacción con dicho entorno.

\begin{figure}[H]
\centering
\includegraphics[width=0.85\textwidth]{slides-images/slide-02.jpg}
\caption{Aprender en entornos dinámicos: el aprendizaje por refuerzo permite que un agente aprenda a tomar decisiones a través de la interacción con el entorno\protect\footnotemark}
\end{figure}
\footnotetext{Diapositiva 2.}

La idea central del aprendizaje por refuerzo profundo difiere esencialmente de los paradigmas de aprendizaje anteriores: ya no depende de un conjunto de datos estático recolectado de antemano, sino que hace que el modelo (Agent) interactúe de manera continua dentro de un \textbf{entorno dinámico} (Environment) —tomando acciones (Action), observando los cambios del entorno (Observation) y recibiendo señales de recompensa (Reward)— para así aprender a tomar decisiones óptimas. Esta forma de aprendizaje tiene un amplio valor de aplicación en ámbitos como el control de robots, las estrategias de juego y la conducción autónoma.

\begin{knowledgebox}{¿Por qué se necesita el aprendizaje por refuerzo?}
En el mundo real, muchas tareas no pueden reducirse simplemente a un problema de aprendizaje supervisado del tipo «dada una entrada, predecir una salida». Por ejemplo, hacer que un robot aprenda a caminar o que una IA aprenda a jugar ajedrez requiere que el agente tome decisiones sucesivas a lo largo de una serie de pasos temporales, y la decisión actual afecta los estados y las recompensas futuros. Este es precisamente el escenario central de aplicación del aprendizaje por refuerzo.
\end{knowledgebox}

\begin{figure}[H]
\centering
\includegraphics[width=0.85\textwidth]{slides-images/slide-03.jpg}
\caption{Ámbitos de aplicación típicos del aprendizaje por refuerzo: robótica y estrategias de juego\protect\footnotemark}
\end{figure}
\footnotetext{Diapositiva 3.}

\section{Comparación de los tres paradigmas de aprendizaje}

Antes de entrar de lleno al aprendizaje por refuerzo, es necesario compararlo sistemáticamente con los dos paradigmas de aprendizaje vistos previamente, para entender su posición particular.

\begin{figure}[H]
\centering
\includegraphics[width=0.95\textwidth]{slides-images/slide-07.jpg}
\caption{Comparación de los tres tipos de problemas de aprendizaje: Supervised Learning, Unsupervised Learning y Reinforcement Learning\protect\footnotemark}
\end{figure}
\footnotetext{Diapositivas, página 7.}

\subsection{Supervised Learning (aprendizaje supervisado)}

\begin{itemize}
    \item \textbf{Datos}: pares $(x, y)$, donde $x$ es el dato de entrada y $y$ es la etiqueta
    \item \textbf{Objetivo}: aprender una función que mapee $x$ a $y$
    \item \textbf{Analogía}: mostrarle al modelo muchas imágenes de manzanas y decirle «esto es una manzana», de modo que el modelo aprenda a reconocer imágenes nuevas de manzanas
\end{itemize}

\subsection{Unsupervised Learning (aprendizaje no supervisado)}

\begin{itemize}
    \item \textbf{Datos}: solo $x$, sin etiquetas
    \item \textbf{Objetivo}: aprender la estructura y la distribución intrínsecas de los datos
    \item \textbf{Analogía}: mostrarle al modelo muchas imágenes de objetos similares (sin decirle qué son), de modo que el modelo aprenda a descubrir sus características comunes
\end{itemize}

\subsection{Reinforcement Learning (aprendizaje por refuerzo)}

\begin{itemize}
    \item \textbf{Datos}: State-Action pairs (pares de estado y acción)
    \item \textbf{Objetivo}: maximizar la \textbf{recompensa acumulada futura} a lo largo de múltiples pasos temporales
    \item \textbf{Analogía}: no es necesario decirle al modelo «esto es una manzana»; en vez de eso, se le deja aprender mediante la interacción que «comer una manzana proporciona nutrientes y así sobrevive»
\end{itemize}

\begin{importantbox}{Diferencia esencial entre los tres paradigmas}
\begin{itemize}
    \item \textbf{Supervised Learning}: aprende el mapeo de entrada a salida a partir de \textbf{datos etiquetados}
    \item \textbf{Unsupervised Learning}: descubre estructuras ocultas a partir de \textbf{datos sin etiquetar}
    \item \textbf{Reinforcement Learning}: aprende una política de decisión óptima mediante la \textbf{interacción con el entorno}, con el objetivo de maximizar el retorno acumulado a largo plazo
\end{itemize}
El algoritmo de aprendizaje (Algorithm) es independiente (agnostic) de la arquitectura de la red (Architecture). Una misma arquitectura (como CNN) puede usarse en distintos paradigmas de aprendizaje.
\end{importantbox}

\subsection{Resumen de la sección}
Los tres paradigmas de aprendizaje corresponden a distintas formas de datos y objetivos de aprendizaje. Lo particular del aprendizaje por refuerzo radica en que: (1) los datos no se recopilan de antemano, sino que se generan dinámicamente mediante la interacción; (2) el objetivo no es predecir etiquetas ni descubrir estructuras, sino aprender una política de conducta óptima que maximice el retorno a largo plazo; (3) la decisión actual afecta los datos que se observarán después (no son i.i.d.).


\section{Conceptos fundamentales del aprendizaje por refuerzo}

Esta sección presenta en detalle los cinco conceptos fundamentales del aprendizaje por refuerzo: Agent (agente), Environment (entorno), Action (acción), Observation/State (observación/estado), Reward (recompensa), así como la relación de interacción entre ellos.

\subsection{Agent y Environment}

\begin{figure}[H]
\centering
\includegraphics[width=0.85\textwidth]{slides-images/slide-11.jpg}
\caption{El circuito de interacción central de RL: el Agent interactúa con el Environment mediante Action y recibe Observation\protect\footnotemark}
\end{figure}
\footnotetext{Diapositiva 11.}

\begin{itemize}
    \item \textbf{Agent} (agente): el sujeto de decisión dentro del sistema, encargado de tomar acciones con base en la observación actual. Por ejemplo: el personaje de Mario en Super Mario, un dron, el controlador de un vehículo autónomo
    \item \textbf{Environment} (entorno): el mundo en el que existe y opera el Agent. Las reglas del entorno definen cómo las acciones del Agent afectan el estado del mundo
\end{itemize}

\subsection{Action y Action Space}

\textbf{Action} (acción) es la operación que el Agent puede ejecutar en el entorno, denotada como $a_t$ (la acción en el instante $t$). El conjunto de todas las acciones posibles se denomina \textbf{Action Space} (espacio de acciones) $\mathcal{A}$.

El espacio de acciones puede ser:
\begin{itemize}
    \item \textbf{Discreto} (Discrete): por ejemplo, las cuatro direcciones arriba, abajo, izquierda y derecha
    \item \textbf{Continuo} (Continuous): por ejemplo, el ángulo del volante (cualquier valor real)
\end{itemize}

\subsection{Observation y State}

Después de que el Agent toma una acción, el Environment devuelve una nueva \textbf{observación} (Observation), que refleja el cambio producido en el entorno por dicha acción. La situación actual percibida por el Agent se denomina \textbf{State} (estado), denotado como $s_t$.

\begin{figure}[H]
\centering
\includegraphics[width=0.85\textwidth]{slides-images/slide-12.jpg}
\caption{Cambio de estado: después de que el Agent actúa, el entorno devuelve el nuevo estado $s_{t+1}$\protect\footnotemark}
\end{figure}
\footnotetext{Diapositiva 12.}

\subsection{Reward y Return}

\textbf{Reward} (recompensa) $r_t$ es la señal escalar que el entorno devuelve al Agent en cada paso de tiempo, y que indica qué tan «bueno» es el estado actual. El objetivo del Agent es maximizar el \textbf{retorno total} (Total Reward / Return):

$$R_t = \sum_{i=t}^{\infty} r_i$$

\begin{itemize}
    \item $R_t$: el retorno total a partir del instante $t$
    \item $r_i$: la recompensa inmediata obtenida en el instante $i$
\end{itemize}

\begin{figure}[H]
\centering
\includegraphics[width=0.85\textwidth]{slides-images/slide-14.jpg}
\caption{El circuito de interacción completo de RL: el Agent envía Action, el Environment devuelve State y Reward\protect\footnotemark}
\end{figure}
\footnotetext{Diapositiva 14.}

\subsection{Retorno descontado (Discounted Return)}

En las aplicaciones reales, no todas las recompensas tienen la misma importancia: las \textbf{recompensas cercanas} suelen valer más que las lejanas (porque el futuro conlleva incertidumbre). Por ello se introduce un \textbf{factor de descuento} $\gamma \in [0, 1)$ para atenuar las recompensas futuras:

$$R_t = \sum_{i=t}^{\infty} \gamma^{i-t} r_i = r_t + \gamma r_{t+1} + \gamma^2 r_{t+2} + \cdots$$

\begin{itemize}
    \item $\gamma$: el factor de descuento (Discount Factor), que controla cuánto peso se le da a las recompensas futuras
    \item Cuanto más cercano a 0 es $\gamma$, más «corto de vista» es el Agent, y solo atiende a la recompensa inmediata
    \item Cuanto más cercano a 1 es $\gamma$, más «largo de vista» es el Agent, y presta más atención al retorno de largo plazo
\end{itemize}

\begin{figure}[H]
\centering
\includegraphics[width=0.85\textwidth]{slides-images/slide-16.jpg}
\caption{Retorno descontado: el factor de descuento $\gamma$ atenúa el peso de las recompensas futuras\protect\footnotemark}
\end{figure}
\footnotetext{Diapositiva 16.}

\begin{warningbox}{¿Por qué se necesita un factor de descuento?}
Si no se usa un factor de descuento ($\gamma = 1$), en un horizonte de tiempo infinito el retorno total puede tender a infinito, lo que vuelve irresoluble el problema de optimización. Además, el factor de descuento también refleja la intuición del mundo real de que «más vale pájaro en mano que ciento volando»: las recompensas futuras conllevan incertidumbre, mientras que la recompensa presente y segura vale más.
\end{warningbox}

\subsection{Q-function (función de valor estado-acción)}

La \textbf{Q-function} es uno de los conceptos más centrales del aprendizaje por refuerzo. Mide el \textbf{retorno total futuro esperado} que el Agent puede obtener al ejecutar la acción $a$ en el estado $s$:

$$Q(s_t, a_t) = \mathbb{E}[R_t \mid s_t, a_t]$$

\begin{figure}[H]
\centering
\includegraphics[width=0.85\textwidth]{slides-images/slide-18.jpg}
\caption{Definición de la Q-function: el retorno total futuro esperado\protect\footnotemark}
\end{figure}
\footnotetext{Diapositiva 18.}

\begin{importantbox}{El significado central de la Q-function}
La pregunta que responde $Q(s, a)$ es: «si tomo la acción $a$ en el estado $s$ y, a partir de ahí, sigo siempre la política óptima, ¿cuánto retorno total puedo obtener?». No solo considera la recompensa inmediata de la acción actual, sino también su efecto sobre todos los \textbf{estados y recompensas futuros}.
\end{importantbox}

\subsection{Policy (política)}

Con la Q-function, el Agent puede establecer una \textbf{política} (Policy), denotada como $\pi(s)$. La política óptima $\pi^*(s)$ consiste en elegir, en cada estado, la acción que maximiza la Q-function:

$$\pi^*(s) = \arg\max_a Q(s, a)$$

\begin{figure}[H]
\centering
\includegraphics[width=0.85\textwidth]{slides-images/slide-19.jpg}
\caption{Cómo usar la Q-function para establecer una política: elegir la acción que maximiza $Q(s,a)$\protect\footnotemark}
\end{figure}
\footnotetext{Diapositiva 19.}

\subsection{Resumen de la sección}

El marco central del aprendizaje por refuerzo puede resumirse como un proceso de interacción en circuito cerrado: el Agent observa el estado actual $s_t$, elige la acción $a_t$ según la política $\pi$, el entorno devuelve el nuevo estado $s_{t+1}$ y la recompensa $r_t$, y el objetivo del Agent es encontrar la política óptima $\pi^*$ que maximice el retorno acumulado descontado. La Q-function es el puente que conecta estado, acción y retorno, mientras que la política es la regla de conducta del Agent.


\section{Los dos grandes paradigmas algorítmicos de Deep RL}

Después de presentar los conceptos básicos de RL, esta sección se centra en cómo implementar algoritmos de aprendizaje por refuerzo usando \textbf{redes neuronales profundas}. Deep RL tiene dos grandes paradigmas:

\begin{figure}[H]
\centering
\includegraphics[width=0.85\textwidth]{slides-images/slide-20.jpg}
\caption{Los dos grandes paradigmas algorítmicos de Deep RL: Value Learning y Policy Learning\protect\footnotemark}
\end{figure}
\footnotetext{Diapositiva 20 de las láminas.}

\begin{importantbox}{Value Learning vs. Policy Learning}
\begin{itemize}
    \item \textbf{Value Learning}(aprendizaje de valor): primero se aprende la Q-function $Q(s, a)$, y luego se determina la acción óptima mediante $a = \arg\max_a Q(s,a)$
    \item \textbf{Policy Learning}(aprendizaje de política): se aprende directamente la función de política $\pi(s)$, y se muestrea la acción a partir de la política $a \sim \pi(s)$
\end{itemize}
Los dos métodos tienen ventajas y desventajas distintas, y son adecuados para diferentes tipos de problemas.
\end{importantbox}

\subsection{Resumen de la sección}
La idea central de Deep RL es usar redes neuronales profundas como aproximador de funciones para aprender la Q-function o la política. Value Learning resuelve el problema de decisión de manera indirecta (primero se aprende $Q$ y luego se decide), mientras que Policy Learning lo resuelve de manera directa (se aprende la política directamente). Las siguientes dos secciones presentarán en detalle, respectivamente, los algoritmos representativos de estos dos paradigmas.
\section{Value Learning: Deep Q-Networks (DQN)}

\subsection{Idea básica de DQN}

Deep Q-Network (DQN) es el algoritmo representativo de Value Learning. Su idea central es: usar una red neuronal profunda para \textbf{aproximar la Q-function}, y después usar la Q-function aprendida para derivar la política óptima.

\begin{figure}[H]
\centering
\includegraphics[width=0.95\textwidth]{slides-images/slide-26.jpg}
\caption{Las dos arquitecturas de DQN: a la izquierda, (state, action) $\to$ Q-value; a la derecha, state $\to$ Q-value de todas las acciones\protect\footnotemark}
\end{figure}
\footnotetext{Slides, página 26.}

DQN tiene dos diseños de arquitectura equivalentes:

\begin{enumerate}
    \item \textbf{Opción uno}: entrada $(s, a)$, salida un único valor $Q(s,a)$. Se necesita llamar a la red una vez por cada acción posible.
    \item \textbf{Opción dos (más eficiente)}: entrada $s$, salida simultánea de los valores Q de todas las acciones $Q(s, a_1), Q(s, a_2), \ldots, Q(s, a_n)$. Basta con una sola pasada hacia adelante.
\end{enumerate}

En la práctica normalmente se usa la opción dos, porque es más eficiente: con una sola llamada a la red se obtiene la evaluación de valor de todas las acciones.

\subsection{Proceso de decisión de DQN}

\begin{figure}[H]
\centering
\includegraphics[width=0.85\textwidth]{slides-images/slide-36.jpg}
\caption{Ejemplo de decisión de DQN: se elige la acción con el valor Q más alto $a_1$ (moverse hacia la izquierda)\protect\footnotemark}
\end{figure}
\footnotetext{Slides, página 36.}

El proceso de decisión es muy directo:
\begin{enumerate}
    \item Se introduce el estado actual $s$ (por ejemplo, la pantalla de un juego de Atari) en DQN
    \item La red genera el valor Q de cada acción posible
    \item Se elige y se ejecuta la acción con el valor Q más grande: $\pi(s) = \arg\max_a Q(s, a)$
\end{enumerate}

\subsection{Ejemplo de Atari Breakout: entender el valor Q}

En clase se usó el juego Atari Breakout para profundizar en el significado de la Q-function. Se consideran dos pares state-action:

\begin{figure}[H]
\centering
\includegraphics[width=0.95\textwidth]{slides-images/slide-22.jpg}
\caption{Comparación de valores Q en Atari Breakout: ¿cuál par (s, a) tiene el valor Q más alto?\protect\footnotemark}
\end{figure}
\footnotetext{Slides, página 22.}

\begin{itemize}
    \item \textbf{Escenario A}: la pelota cae en línea recta hacia el centro de la paleta, y la paleta permanece inmóvil
    \item \textbf{Escenario B}: la pelota vuela hacia un costado y la paleta se mueve a la derecha para alcanzarla
\end{itemize}

\begin{knowledgebox}{Intuición sobre el valor Q}
Intuitivamente el valor Q del escenario B podría parecer más alto (la pelota rebota hacia el costado y golpea más ladrillos), pero los experimentos reales muestran que la estrategia del escenario B (el golpe lateral), aunque puede dar una puntuación mayor en una sola jugada, deja la paleta desplazada hacia el costado y hace que la velocidad de devolución sea más lenta, lo que a largo plazo resulta desfavorable. Esto muestra que el valor Q mide el \textbf{retorno acumulado a largo plazo} y no la recompensa de un solo paso; la intuición humana puede no estimar con precisión el valor Q, y precisamente por eso se necesita una red profunda para aprenderlo.
\end{knowledgebox}

\subsection{Entrenamiento de DQN}

El núcleo del entrenamiento de DQN es definir un valor objetivo y una función de pérdida adecuados.

\textbf{Target} (valor objetivo): se construye el objetivo de entrenamiento a partir de la ecuación de Bellman:

$$\text{target} = r + \gamma \max_{a'} Q(s', a')$$

\begin{itemize}
    \item $r$: la recompensa inmediata obtenida en el paso actual
    \item $s'$: el nuevo estado al que se llega tras ejecutar la acción
    \item $\max_{a'} Q(s', a')$: el valor Q más grande entre todas las acciones posibles en el nuevo estado (suponiendo que después se sigue la política óptima)
\end{itemize}

\textbf{Q-Loss} (función de pérdida):

$$\mathcal{L} = \mathbb{E}\left[\left\| \underbrace{\left(r + \gamma \max_{a'} Q(s', a')\right)}_{\text{target}} - \underbrace{Q(s, a)}_{\text{predicted}} \right\|^2\right]$$

\begin{figure}[H]
\centering
\includegraphics[width=0.95\textwidth]{slides-images/slide-30.jpg}
\caption{Entrenamiento de DQN: la Q-Loss se define como el error cuadrático medio entre el valor Q target y el predicted\protect\footnotemark}
\end{figure}
\footnotetext{Slides, página 30.}

\begin{importantbox}{Intuición central del entrenamiento de DQN}
El entrenamiento de DQN consiste, en esencia, en hacer que el valor Q predicho por la red se acerque gradualmente al valor target dado por la ecuación de Bellman. El target está formado por «la recompensa inmediata más el retorno futuro óptimo descontado», con lo que se logra el objetivo de «usar el retorno futuro para guiar la decisión actual». El proceso de entrenamiento realiza la retropropagación y la actualización de parámetros minimizando el error cuadrático medio entre el Q predicted y el Q target.
\end{importantbox}

\subsection{Resultados de DQN en juegos de Atari}

Google DeepMind publicó en 2015 en Nature el artículo histórico de DQN, en el que se usó la misma arquitectura de red y los mismos hiperparámetros para hacer pruebas en 49 juegos de Atari.

\begin{figure}[H]
\centering
\includegraphics[width=0.95\textwidth]{slides-images/slide-32.jpg}
\caption{Arquitectura de red de DQN: las capas convolucionales procesan la pantalla del juego y las capas totalmente conectadas generan el valor Q de cada acción\protect\footnotemark}
\end{figure}
\footnotetext{Slides, página 32. Mnih y colaboradores, Nature 2015.}

La arquitectura de red incluye:
\begin{itemize}
    \item Varias capas de \textbf{Convolution} (capas convolucionales): procesan los píxeles de la pantalla del juego sin procesar
    \item Varias capas \textbf{Fully Connected} (totalmente conectadas): mapean las características al valor Q de cada acción
    \item Capa de salida: cada nodo corresponde a una acción posible (por ejemplo, cada combinación de direcciones y botones del control del juego)
\end{itemize}

\begin{figure}[H]
\centering
\includegraphics[width=0.95\textwidth]{slides-images/slide-33.jpg}
\caption{Desempeño de DQN en juegos de Atari: más del 50\% de los juegos alcanzan o superan el nivel humano\protect\footnotemark}
\end{figure}
\footnotetext{Slides, página 33. Mnih y colaboradores, Nature 2015.}

\begin{knowledgebox}{Importancia histórica de DQN}
La importancia de DQN radica en que demostró que un sistema de aprendizaje por refuerzo profundo \textbf{de extremo a extremo} puede aprender directamente, a partir de píxeles sin procesar, políticas complejas. Antes de esto, el RL solía requerir el diseño manual de características (Feature Engineering). DQN abrió una nueva era del Deep RL y fue una de las tecnologías clave que llevaron a que Google adquiriera DeepMind.
\end{knowledgebox}

\subsection{Limitaciones de Q-Learning}

\begin{figure}[H]
\centering
\includegraphics[width=0.85\textwidth]{slides-images/slide-34.jpg}
\caption{Dos grandes limitaciones de Q-Learning: no puede manejar espacios de acción continuos y solo puede aprender políticas deterministas\protect\footnotemark}
\end{figure}
\footnotetext{Slides, página 34.}

Aunque DQN logró un enorme éxito, el método Q-Learning tiene dos limitaciones importantes:

\textbf{Problema de complejidad} (Complexity):
\begin{itemize}
    \item Solo puede manejar espacios de acción \textbf{discretos y finitos}
    \item Para espacios de acción continuos (como controlar el ángulo del volante), no es posible enumerar todas las acciones posibles
\end{itemize}

\textbf{Problema de flexibilidad} (Flexibility):
\begin{itemize}
    \item La política se deriva de la Q-function de manera determinista mediante $\arg\max$
    \item No puede aprender una \textbf{política estocástica} (Stochastic Policy), y en algunos escenarios la aleatoriedad es beneficiosa
\end{itemize}

\begin{warningbox}{¿Cuándo no es aplicable Q-Learning?}
Cuando el espacio de acción es continuo (como los ángulos de las articulaciones de un robot, o la fuerza del acelerador y el freno de un automóvil), Q-Learning no puede aplicarse directamente, porque necesita calcular el valor Q para \textbf{todas las acciones posibles} y tomar el máximo. En un espacio continuo hay infinitas acciones, lo cual es computacionalmente inviable. En ese caso hay que recurrir a los métodos de Policy Gradient.
\end{warningbox}

\subsection{Resumen de la sección}

DQN aproxima la Q-function mediante una red neuronal profunda, con lo que logra un aprendizaje de extremo a extremo desde los píxeles sin procesar hasta la decisión. Su entrenamiento se basa en construir el target con la ecuación de Bellman y actualizar los parámetros de la red minimizando la Q-Loss. DQN alcanzó resultados que superan el nivel humano en juegos de Atari, pero, al estar limitado a espacios de acción discretos y políticas deterministas, no puede manejar directamente problemas más complejos como el control continuo.

\section{Policy Learning: método de Policy Gradient}

Para superar las limitaciones de Q-Learning, esta sección presenta otra clase de algoritmos de Deep RL—\textbf{Policy Gradient} (gradiente de política), que aprende directamente la función de política $\pi(s)$, sin necesidad de aprender la Q-function intermedia.

\subsection{Idea central de Policy Gradient}

\begin{figure}[H]
\centering
\includegraphics[width=0.85\textwidth]{slides-images/slide-38.jpg}
\caption{Política en un espacio de acciones discreto: salida de la distribución de probabilidad de cada acción $P(a|s)$\protect\footnotemark}
\end{figure}
\footnotetext{Slides, página 38.}

A diferencia de DQN, el método Policy Gradient ya no estima el valor Q de cada acción, sino que produce directamente una \textbf{distribución de probabilidad sobre las acciones}. La entrada de la red es el estado $s$, y la salida es $P(a|s)$—la probabilidad de tomar cada acción en ese estado.

Para un espacio de acciones \textbf{discreto}, la salida es una distribución de probabilidad softmax (similar a una tarea de clasificación):

$$P(a = \text{left} | s) = 0.6, \quad P(a = \text{stay} | s) = 0.25, \quad P(a = \text{right} | s) = 0.15$$

\subsection{Manejo de espacios de acciones continuos}

La mayor ventaja de Policy Gradient es que puede manejar de forma natural \textbf{espacios de acciones continuos}.

\begin{figure}[H]
\centering
\includegraphics[width=0.95\textwidth]{slides-images/slide-40.jpg}
\caption{Policy Gradient en espacios de acciones continuos: la red produce los parámetros de una distribución gaussiana $(\mu, \sigma^2)$\protect\footnotemark}
\end{figure}
\footnotetext{Slides, página 40.}

Para un espacio de acciones continuo, la red ya no produce la probabilidad de cada acción, sino que produce los \textbf{parámetros de una distribución de probabilidad}. Por ejemplo, supongamos que la acción sigue una distribución gaussiana:

$$P(a|s) = \mathcal{N}(\mu, \sigma^2)$$

La red produce la media $\mu$ y la varianza $\sigma^2$, y luego se muestrea de esa distribución para obtener una acción concreta:

$$\pi(s) \sim P(a|s) = \mathcal{N}(\mu, \sigma^2)$$

\begin{importantbox}{Ventajas clave de Policy Gradient}
\begin{itemize}
    \item \textbf{Manejo de espacios de acciones continuos}: al producir los parámetros de una distribución de probabilidad (como $\mu$ y $\sigma^2$ de una distribución gaussiana), puede manejar de forma natural problemas de control continuo
    \item \textbf{Aprendizaje de una política estocástica}: la salida es una distribución de probabilidad en lugar de una acción determinista, lo que da soporte natural a Stochastic Policy
    \item \textbf{Optimización directa del objetivo}: optimiza directamente la política para maximizar el retorno esperado, sin necesidad de aprender la Q-function intermedia
\end{itemize}
\end{importantbox}

\subsection{Algoritmo de entrenamiento de Policy Gradient}

\begin{figure}[H]
\centering
\includegraphics[width=0.85\textwidth]{slides-images/slide-44.jpg}
\caption{Flujo de entrenamiento de Policy Gradient: algoritmo de entrenamiento de cinco pasos\protect\footnotemark}
\end{figure}
\footnotetext{Slides, página 44.}

El entrenamiento de Policy Gradient sigue el siguiente flujo de cinco pasos:

\begin{enumerate}
    \item \textbf{Inicializar el Agent}: inicializar de forma aleatoria los parámetros de la red de política
    \item \textbf{Ejecutar la política hasta la terminación}: hacer que el Agent ejecute un episode completo en el entorno siguiendo la política actual
    \item \textbf{Registrar todos los (state, action, reward)}: recopilar todos los datos de trayectoria de todo el episode
    \item \textbf{Reducir la probabilidad de las acciones de baja recompensa}: para las acciones que llevan a un retorno bajo, reducir la probabilidad de que sean elegidas
    \item \textbf{Aumentar la probabilidad de las acciones de alta recompensa}: para las acciones que llevan a un retorno alto, aumentar la probabilidad de que sean elegidas
\end{enumerate}

\subsection{Función de pérdida de Policy Gradient}

El núcleo matemático del entrenamiento consiste en definir una función de pérdida adecuada:

$$\text{loss} = -\log P(a_t | s_t) \cdot R_t$$

\begin{figure}[H]
\centering
\includegraphics[width=0.95\textwidth]{slides-images/slide-46.jpg}
\caption{Función de pérdida y fórmula de actualización del gradiente de Policy Gradient\protect\footnotemark}
\end{figure}
\footnotetext{Slides, página 46. Williams Reinforcement Learning 1992.}

\begin{itemize}
    \item $-\log P(a_t | s_t)$: el logaritmo negativo de la verosimilitud (log-likelihood) de elegir la acción $a_t$
    \item $R_t$: el retorno acumulado descontado a partir del instante $t$
\end{itemize}

La regla de actualización por descenso de gradiente es:

$$w' = w - \nabla \text{loss} = w + \nabla \log P(a_t | s_t) \cdot R_t$$

\begin{importantbox}{Comprensión intuitiva de la función de pérdida de Policy Gradient}
El diseño de esta función de pérdida es muy ingenioso:
\begin{itemize}
    \item Cuando $R_t > 0$ (retorno alto), el signo negativo del loss hace que la actualización del gradiente \textbf{aumente} $P(a_t|s_t)$, es decir, incremente la probabilidad de esa acción
    \item Cuando $R_t < 0$ (retorno bajo o penalización), la actualización del gradiente \textbf{reduce} $P(a_t|s_t)$, es decir, disminuye la probabilidad de esa acción
    \item El valor absoluto de $R_t$ determina la magnitud de la actualización—cuanto mayor es el retorno, más fuerte es el ajuste de la política
\end{itemize}
Esta es la idea central del algoritmo REINFORCE (Williams, 1992).
\end{importantbox}

\begin{warningbox}{El problema de alta varianza de Policy Gradient}
Un problema conocido del método Policy Gradient es la \textbf{alta varianza} (High Variance). Como se usa el muestreo de Monte Carlo para estimar el retorno, la fluctuación del retorno entre distintos episode puede ser muy grande, lo que genera un ruido considerable en la estimación del gradiente e inestabilidad en el entrenamiento. Los métodos habituales para mitigar esto incluyen: usar un baseline para reducir la varianza (como en los métodos Actor-Critic), aumentar el número de muestras, y usar reward normalization, entre otros.
\end{warningbox}

\subsection{Resumen de la sección}

El método Policy Gradient aprende directamente la función de política y optimiza la política mediante «aumentar la probabilidad de las acciones de alto retorno y reducir la probabilidad de las acciones de bajo retorno». Su mayor ventaja es que puede manejar espacios de acciones continuos y aprender políticas estocásticas. La función de pérdida se construye a partir del producto entre el logaritmo de la verosimilitud de la acción y el retorno descontado, y el retorno esperado se optimiza mediante ascenso de gradiente.

\section{Aplicaciones del aprendizaje por refuerzo en el mundo real}

Esta sección presenta tres aplicaciones importantes del aprendizaje por refuerzo profundo: el RL de extremo a extremo en la conducción autónoma, AlphaGo y AlphaZero en el go, y, en años recientes, el RLHF en la alineación de los LLM.

\subsection{Conducción autónoma: de la simulación al despliegue real}

\begin{figure}[H]
\centering
\includegraphics[width=0.95\textwidth]{slides-images/slide-48.jpg}
\caption{VISTA: simulador de conducción autónoma de alta fidelidad desarrollado por el MIT\protect\footnotemark}
\end{figure}
\footnotetext{Diapositiva 48. Amini+ IEEE RA-L 2020.}

El profesor Amini presentó su trabajo de investigación en el MIT, en el que usa el método de gradiente de política para entrenar sistemas de conducción autónoma. Las innovaciones clave de este trabajo incluyen:

\textbf{El simulador VISTA}: un entorno de simulación de alta fidelidad basado en datos reales de conducción. A diferencia de la simulación tradicional por gráficos computacionales, VISTA usa video real de conducción para sintetizar nuevos puntos de vista, lo que garantiza que las imágenes simuladas tengan un realismo fotográfico.

\textbf{Entrenamiento de RL de extremo a extremo}:
\begin{itemize}
    \item El estado del agente es la imagen de la cámara de conducción
    \item La acción es el ángulo continuo del volante
    \item El diseño de la recompensa es extremadamente disperso: solo se otorga una penalización cuando el vehículo se sale del carril o sufre una colisión
    \item La red de política produce los parámetros de una distribución gaussiana, de la cual se muestrea el ángulo de dirección concreto
\end{itemize}

\begin{figure}[H]
\centering
\includegraphics[width=0.95\textwidth]{slides-images/slide-49.jpg}
\caption{El primer automóvil autónomo del mundo entrenado completamente mediante RL en simulación y desplegado directamente en carreteras reales\protect\footnotemark}
\end{figure}
\footnotetext{Diapositiva 49. Amini+ IEEE RA-L 2020.}

\begin{knowledgebox}{De la simulación a la realidad (Sim-to-Real Transfer)}
Lo innovador de este trabajo es que el modelo se entrena \textbf{por completo en un entorno simulado}, sin haber tenido nunca contacto con datos reales de carretera, pero puede \textbf{desplegarse directamente en un automóvil real} y conducir con seguridad. Este es el primer automóvil autónomo de tamaño real del mundo entrenado por completo mediante RL en simulación y desplegado en el mundo real. Lo esencial está en que la alta fidelidad del simulador VISTA cierra la brecha entre la simulación y la realidad (Domain Gap).
\end{knowledgebox}

\begin{figure}[H]
\centering
\includegraphics[width=0.95\textwidth]{slides-images/slide-47.jpg}
\caption{La dificultad de aplicar el entrenamiento de RL en el mundo real: ejecutar la política (paso 2) puede provocar peligros\protect\footnotemark}
\end{figure}
\footnotetext{Diapositiva 47.}

\begin{warningbox}{El riesgo de entrenar RL en entornos reales}
El paso 2 del algoritmo de entrenamiento por gradiente de política, "ejecutar la política hasta la terminación", puede ser sumamente peligroso en un entorno real. En el caso de la conducción autónoma, "terminación" puede significar una colisión o un accidente. Por ello, el \textbf{entorno simulado} es el requisito previo para que el RL se aplique en dominios críticos para la seguridad: el agente puede cometer errores y aprender con seguridad en la simulación, sin causar daño real.
\end{warningbox}

\subsection{El go: AlphaGo y AlphaZero}

El go (Go) es una de las aplicaciones más llamativas del RL profundo. La complejidad del go supera con mucho a la del ajedrez: el número de posiciones legales en un tablero de 19x19 es de aproximadamente $2.08 \times 10^{170}$, más que la cantidad de átomos del universo.

\begin{figure}[H]
\centering
\includegraphics[width=0.95\textwidth]{slides-images/slide-52.jpg}
\caption{El espacio de búsqueda del go: el número de posiciones legales supera la cantidad de átomos del universo\protect\footnotemark}
\end{figure}
\footnotetext{Diapositiva 52. Fuente: Wikipedia.}

\textbf{AlphaGo} (Silver+ Science 2016) empleó un proceso de entrenamiento en varias etapas:

\begin{figure}[H]
\centering
\includegraphics[width=0.95\textwidth]{slides-images/slide-54.jpg}
\caption{El proceso de entrenamiento de AlphaGo: de los datos de expertos humanos al autojuego\protect\footnotemark}
\end{figure}
\footnotetext{Diapositiva 54. Silver+ Science 2016.}

\begin{enumerate}
    \item \textbf{Preentrenamiento supervisado}: usar partidas de expertos humanos para entrenar la red de política (Policy Network) como una tarea de clasificación (predecir la jugada de un jugador humano)
    \item \textbf{Autojuego con aprendizaje por refuerzo}: usar el método de gradiente de política del RL para que la red de política siga mejorando mediante el autojuego, hasta alcanzar un nivel \textbf{superior al humano}
    \item \textbf{Entrenamiento de la red de valor}: usar los datos generados por el autojuego para entrenar la red de valor (Value Network), que evalúa la probabilidad de victoria de una posición (tarea de regresión)
\end{enumerate}

En 2016, AlphaGo venció 4 a 1 al campeón mundial de go Lee Sedol, un hito en la historia de la IA.

\textbf{AlphaZero} (Silver+ Science 2018) fue un paso más allá: \textbf{abandonó por completo los datos humanos} y se entrenó desde cero únicamente mediante aprendizaje por refuerzo con autojuego:

\begin{figure}[H]
\centering
\includegraphics[width=0.85\textwidth]{slides-images/slide-57.jpg}
\caption{AlphaZero: aprendizaje desde cero por completo mediante autojuego, con la puntuación Elo subiendo rápidamente conforme avanzan los pasos de entrenamiento\protect\footnotemark}
\end{figure}
\footnotetext{Diapositiva 57. Silver+ Science 2018.}

\begin{importantbox}{El salto de AlphaGo a AlphaZero}
\begin{itemize}
    \item \textbf{AlphaGo}: necesita partidas de expertos humanos como punto de partida (preentrenamiento supervisado + ajuste fino con RL)
    \item \textbf{AlphaZero}: comienza por completo desde cero (tabula rasa), sin usar ningún dato humano, y alcanza, solo mediante \textbf{autojuego más aprendizaje por refuerzo}, un nivel superior al de todas las versiones anteriores
    \item AlphaZero demostró que, en juegos de información perfecta, el autojuego con RL puro puede \textbf{descubrir estrategias que los humanos nunca imaginaron}, superando los límites de la cognición humana
\end{itemize}
\end{importantbox}

\subsection{Resumen de la sección}

El RL profundo ha logrado resultados innovadores en ámbitos como la conducción autónoma y el go. El RL de extremo a extremo en la conducción autónoma muestra la capacidad del gradiente de política para abordar tareas de control continuo, y el simulador es la clave para garantizar un entrenamiento seguro. La serie AlphaGo y AlphaZero demuestra que el RL profundo combinado con el autojuego puede alcanzar, e incluso superar, el nivel humano en problemas de decisión extremadamente complejos.

\section{Retos y fronteras del RL: diseño de la función de recompensa}

En clase, al discutir las aplicaciones prácticas de DQN y Policy Gradient, se enfatizó en particular un problema ampliamente subestimado en el RL pero extremadamente crítico------\textbf{el diseño de la función de recompensa} (Reward Design / Reward Engineering).

\subsection{El origen de la recompensa}

En los juegos de Atari, la recompensa está definida de forma natural por las reglas del juego (por ejemplo, romper un ladrillo suma puntos, perder una vida los resta), y quien diseña el sistema no necesita ocuparse de ello. Pero en las aplicaciones del mundo real, la definición de la recompensa está lejos de ser evidente:

\begin{itemize}
    \item En la conducción autónoma, ¿qué es "conducir bien"? ¿Seguridad? ¿Eficiencia? ¿Comodidad?
    \item En el control robótico, ¿cómo se cuantifica la calidad con la que se completa una tarea?
    \item En los modelos de lenguaje, ¿qué es "una buena respuesta"?
\end{itemize}

\begin{warningbox}{Consecuencias de un diseño de recompensa deficiente (Reward Hacking)}
Si la función de recompensa está mal diseñada, el Agent puede encontrar formas de \textbf{explotar fallas de la función de recompensa} para obtener una recompensa alta sin haber completado realmente la tarea prevista. Por ejemplo:
\begin{itemize}
    \item Un robot de limpieza descubre que "esconder la basura" obtiene una recompensa mayor que "limpiar de verdad"
    \item Una IA de carreras descubre que "dar vueltas repetidamente para recolectar objetos de aceleración" puntúa más que "terminar la carrera"
\end{itemize}
Esto se conoce como Reward Hacking o Reward Misspecification, y es una de las trampas más comunes en las aplicaciones de RL.
\end{warningbox}

\subsection{Recompensa dispersa y recompensa densa}

El grado de dispersión de la señal de recompensa afecta directamente la eficiencia del aprendizaje:

\begin{itemize}
    \item \textbf{Recompensa densa} (Dense Reward): hay retroalimentación clara en cada paso (como en Atari, donde cada ladrillo roto otorga puntos), y el aprendizaje es más rápido
    \item \textbf{Recompensa dispersa} (Sparse Reward): solo hay retroalimentación al final, en caso de éxito o fracaso (como en la conducción autónoma, donde solo hay penalización en caso de colisión); el aprendizaje es difícil, pero el diseño de la recompensa es más simple
\end{itemize}

\begin{knowledgebox}{RLHF: usar la retroalimentación humana para definir la recompensa}
Para tareas como las de los modelos de lenguaje, es difícil definir "una buena salida" con una fórmula matemática simple. \textbf{RLHF} (Reinforcement Learning from Human Feedback) entrena un Reward Model que aproxima las preferencias humanas al hacer que personas ordenen y califiquen las salidas del modelo, y después usa ese Reward Model para guiar la optimización de la política. Esta es una de las técnicas centrales del alineamiento (Alignment) de los LLM modernos, como ChatGPT, y en esencia es Policy Gradient combinado con un modelo de recompensa definido por seres humanos.
\end{knowledgebox}

\subsection{Resumen de la sección}

La función de recompensa es el alma del sistema de RL------define el objetivo de optimización del Agent. Un diseño de recompensa demasiado simple puede provocar Reward Hacking, mientras que uno demasiado complejo es difícil de llevar a la ingeniería. RLHF ofrece un método para definir la recompensa mediante el juicio humano en tareas difíciles de formalizar, y es una de las rutas técnicas importantes del alignment de la IA actual.
\section{Síntesis y ampliación}

\subsection{Resumen de la clase}

\begin{figure}[H]
\centering
\includegraphics[width=0.95\textwidth]{slides-images/slide-58.jpg}
\caption{Resumen de Deep RL: tres bloques -- conceptos básicos, Q-Learning y Policy Gradient\protect\footnotemark}
\end{figure}
\footnotetext{Diapositivas, página 58.}

Esta clase presentó de manera sistemática el contenido central del Deep Reinforcement Learning, dividido principalmente en tres partes:

\textbf{Uno. Conceptos básicos de RL}
\begin{itemize}
    \item El Agent interactúa con el Environment mediante Action
    \item Los pares State-Action constituyen los datos básicos del aprendizaje
    \item El objetivo es maximizar el retorno acumulado con descuento $R_t = \sum \gamma^i r_i$
    \item La Q-function $Q(s,a)$ mide el retorno futuro esperado de un par estado-acción
\end{itemize}

\textbf{Dos. Value Learning / Q-Learning}
\begin{itemize}
    \item Se usa DQN para aproximar la Q-function
    \item La política se deriva de forma determinista mediante $\arg\max_a Q(s,a)$
    \item Aplicable a espacios de acción discretos
    \item Logró un desempeño superior al humano en juegos de Atari
\end{itemize}

\textbf{Tres. Policy Gradient}
\begin{itemize}
    \item Aprende y optimiza directamente la función de política $\pi(s)$
    \item Aplicable a espacios de acción continuos
    \item Se entrena mediante $\text{loss} = -\log P(a_t|s_t) \cdot R_t$
    \item Logró avances en aplicaciones complejas como la conducción autónoma y el Go
\end{itemize}

\subsection{Lecturas adicionales}

\begin{itemize}
    \item \textbf{Artículo original de DQN}: Mnih et al., ``Human-level control through deep reinforcement learning,'' \textit{Nature}, 2015
    \item \textbf{Algoritmo REINFORCE}: Williams, ``Simple statistical gradient-following algorithms for connectionist reinforcement learning,'' \textit{Machine Learning}, 1992
    \item \textbf{AlphaGo}: Silver et al., ``Mastering the game of Go with deep neural networks and tree search,'' \textit{Nature}, 2016
    \item \textbf{AlphaZero}: Silver et al., ``A general reinforcement learning algorithm that masters chess, shogi, and Go through self-play,'' \textit{Science}, 2018
    \item \textbf{Simulador VISTA}: Amini et al., ``Learning Robust Control Policies for End-to-End Autonomous Driving From Data-Driven Simulation,'' \textit{IEEE RA-L}, 2020
    \item \textbf{RLHF}: Ouyang et al., ``Training language models to follow instructions with human feedback,'' \textit{NeurIPS}, 2022
    \item \textbf{PPO}: Schulman et al., ``Proximal Policy Optimization Algorithms,'' 2017------la variante de Policy Gradient más popular actualmente
    \item \textbf{Métodos Actor-Critic}: combinan las ventajas de Value Learning y Policy Gradient, aprendiendo simultáneamente la función de valor y la función de política
    \item \textbf{Sitio del curso}: \url{https://introtodeeplearning.com}
\end{itemize}

\begin{knowledgebox}{El futuro del Deep RL}
El aprendizaje por refuerzo profundo sigue siendo un campo en rápido desarrollo. Los focos de investigación actuales incluyen:
\begin{itemize}
    \item \textbf{Sample Efficiency}: cómo alcanzar el mismo desempeño con menos datos de interacción
    \item \textbf{Multi-Agent RL}: escenarios en los que varios Agent aprenden y colaboran o compiten simultáneamente
    \item \textbf{Offline RL}: entrenamiento sin conexión a partir de datos históricos existentes, evitando el costo y el riesgo de la interacción en línea
    \item \textbf{Foundation Models + RL}: incorporar el conocimiento de los modelos preentrenados a gran escala al RL para acelerar el aprendizaje
    \item \textbf{RLHF y AI Alignment}: entrenar sistemas de IA más seguros y más útiles mediante retroalimentación humana
\end{itemize}
\end{knowledgebox}


%% --- Fin del contenido --- %%

\end{document}
